\section*{Capítulo \ref{ch:b1:intermolecular-forces-solvents} --- Fuerzas intermoleculares y disolventes}

\begin{solution}{exo:b1:intermolecular-forces-solvents:1}
Argón: solo London. Cloruro de hidrógeno: Keesom, Debye y London, con
predominio de London (\cref{exo:b1:intermolecular-forces-solvents:8}); el
cloro es demasiado grande para formar \omterm{def:b1:intermolecular-forces-solvents:hbond}{enlaces de hidrógeno} fuertes. Agua:
dominan los \omterm{def:b1:intermolecular-forces-solvents:hbond}{enlaces de hidrógeno}, junto con los tres términos de van der
Waals. Metano: solo London. Diyodo: solo London, intensa porque la
molécula es grande y muy polarizable.
\end{solution}

\begin{solution}{exo:b1:intermolecular-forces-solvents:2}
Entre sus propias moléculas: \ce{CH3OH} (\ce{O-H}), \ce{CH3NH2}
(\ce{N-H}) y \ce{HF}. \ce{CH3OCH3} y \ce{CH3F} tienen \omterm{def:b1:lewis-resonance-vsepr:lewis}{pares no enlazantes}
en O o en F, pero ningún hidrógeno unido a ellos: solo pueden aceptar
\omterm{def:b1:intermolecular-forces-solvents:hbond}{enlaces de hidrógeno}, del agua por ejemplo. \ce{CH4} ni da ni acepta.
\end{solution}

\begin{solution}{exo:b1:intermolecular-forces-solvents:3}
Agua y etanol: polares, próticos. Propanona y acetonitrilo: polares,
apróticos. Diclorometano: poco polar, aprótico. Hexano: apolar, aprótico.
\end{solution}

\begin{solution}{exo:b1:intermolecular-forces-solvents:4}
El diyodo es apolar: en el hexano cambia \omterm{def:b1:intermolecular-forces-solvents:vdw}{interacciones de London} por
\omterm{def:b1:intermolecular-forces-solvents:vdw}{interacciones de London}, mientras que en el agua rompería \omterm{def:b1:intermolecular-forces-solvents:hbond}{enlaces de
hidrógeno} sin reemplazarlos. El cloruro de sodio es iónico: solo un
disolvente de permitividad alta puede separar sus iones, y solo uno polar
puede solvatarlos; el hexano ($\varepsilon_r = 1.89$) no hace ninguna de
las dos cosas.
\end{solution}

\begin{solution}{exo:b1:intermolecular-forces-solvents:5}
El metoximetano no tiene \ce{O-H}: no puede dar \omterm{def:b1:intermolecular-forces-solvents:hbond}{enlaces de hidrógeno} y es
el que hierve a menor temperatura. El metanol y el etanol forman ambos
\omterm{def:b1:intermolecular-forces-solvents:hbond}{enlaces de hidrógeno}; el etanol, con un \ce{CH2} más, tiene \omterm{def:b1:intermolecular-forces-solvents:vdw}{interacciones
de London} más intensas. El ácido etanoico forma dímeros unidos por
\omterm{def:b1:intermolecular-forces-solvents:hbond}{enlaces de hidrógeno} y es el más pesado: es el que hierve a mayor
temperatura.
\end{solution}

\begin{solution}{exo:b1:intermolecular-forces-solvents:6}
En el vacío, $F = e^2/(4\pi\varepsilon_0 r^2) = (1.602 \times
10^{-19})^2/(4\pi \times 8.854 \times 10^{-12} \times (5.00 \times
10^{-10})^2) = \qty{9.23e-10}{N}$. En agua, $F/78.4 = \qty{1.18e-11}{N}$;
en hexano, $F/1.89 = \qty{4.88e-10}{N}$. El agua debilita la atracción
41 veces más que el hexano: puede mantener separados los iones; el
hexano, no.
\end{solution}

\begin{solution}{exo:b1:intermolecular-forces-solvents:7}
Mezclar la propanona con agua rompe algunos \omterm{def:b1:intermolecular-forces-solvents:hbond}{enlaces de hidrógeno}
agua--agua, pero forma otros nuevos entre el agua y el oxígeno del
\ce{C=O}, además de interacciones dipolo--dipolo: el balance es
favorable. El hexano solo ofrece al agua débiles \omterm{def:b1:intermolecular-forces-solvents:vdw}{interacciones de London}
a cambio de los \omterm{def:b1:intermolecular-forces-solvents:hbond}{enlaces de hidrógeno} que esta rompería: los dos líquidos
permanecen separados.
\end{solution}

\begin{solution}{exo:b1:intermolecular-forces-solvents:8}
\ce{HBr} tiene más electrones y una nube más grande y polarizable: su
\omterm{def:b1:intermolecular-forces-solvents:vdw}{interacción de London} supera a la de \ce{HCl} en más de lo que se queda
corta su \omterm{def:b1:intermolecular-forces-solvents:vdw}{interacción de Keesom}. Domina London.
\end{solution}

\begin{solution}{exo:b1:intermolecular-forces-solvents:9}
\chemfig{CH_3-C(=[:60]O)-[:-60]O-[:0]H} frente a su pareja, con dos
enlaces $\ce{O-H}\cdots\ce{O=C}$ que cierran un anillo de ocho miembros.
En el benceno, apolar y aprótico, nada compite con esos enlaces: el ácido
está dimerizado y se comporta como una partícula de unos
\qty{120}{g/mol}. En el agua, cada molécula de ácido forma en su lugar
\omterm{def:b1:intermolecular-forces-solvents:hbond}{enlaces de hidrógeno} con el agua: los dímeros se rompen.
\end{solution}

\begin{solution}{exo:b1:intermolecular-forces-solvents:10}
El cloruro de sodio está formado por iones: el agua, de permitividad
alta, los separa (disociación) y los hidrata (el oxígeno hacia
\ce{Na+}, el hidrógeno hacia \ce{Cl-}). El cloruro de hidrógeno es una
\omterm{def:b1:lewis-resonance-vsepr:dipole}{molécula polar}: el agua, polar, convierte el enlace \ce{H-Cl} en un par
de iones \ce{H3O+ Cl-} (ionización), separa el par (disociación) e
hidrata los iones. El hexano es apolar y de baja permitividad: ni ioniza
ni disocia, de modo que no hay iones que transporten la corriente.
\end{solution}

\begin{solution}{exo:b1:intermolecular-forces-solvents:11}
$(174.1 - 36.1)/5 = \qty{27.6}{\celsius}$ por \ce{CH2}. El último
incremento (del nonano al decano) es de \qty{23.4}{\celsius}, de modo que
el undecano debería hervir hacia $174 + 22 \approx \qty{196}{\celsius}$.
Cada \ce{CH2} añade la misma contribución de London, pero una fracción
cada vez menor del total de la molécula, mientras que la temperatura
necesaria para vencer las interacciones ya es alta.
% ledger: bp:C9H20
\end{solution}

\begin{solution}{exo:b1:intermolecular-forces-solvents:12}
Parte \omterm{def:b1:intermolecular-forces-solvents:amphiphile}{hidrófila}: la cabeza sulfato \ce{-OSO3^-} con su contraión
\ce{Na+}; parte \omterm{def:b1:intermolecular-forces-solvents:amphiphile}{hidrófoba}: la cadena de doce carbonos. En la superficie,
las cabezas en el agua y las colas en el aire; en una micela, las colas
dentro y las cabezas fuera. Las colas se disuelven en la grasa de la
mancha, las cabezas la mantienen en contacto con el agua, y la grasa se
desprende en forma de micelas.
\end{solution}

\begin{solution}{pb:b1:intermolecular-forces-solvents:1}
\textbf{1.} La \omterm{def:b1:intermolecular-forces-solvents:vdw}{interacción de London}, la única entre átomos apolares.
\textbf{2.} Del helio al xenón crecen el número de electrones y el tamaño
y, con ellos, la \omterm{def:b1:periodicity:polarisability}{polarizabilidad} y la atracción de London.
\textbf{3.} El helio tiene dos electrones fuertemente retenidos por su
núcleo: es el átomo menos polarizable, con la atracción de London más
débil.
\textbf{4.} $68.8 - 36.1 = \qty{32.7}{\celsius}$: un \ce{CH2} aporta
electrones y superficie de contacto y, por tanto, atracción de London.
\textbf{5.} London, que crece con la \omterm{def:b1:periodicity:polarisability}{polarizabilidad}: $\ce{CH4} < \ce{SiH4} <
\ce{GeH4}$.
\textbf{6.} El carbono no es lo bastante electronegativo para polarizar
mucho el enlace \ce{C-H}, y el metano no tiene ningún \omterm{def:b1:lewis-resonance-vsepr:lewis}{par no enlazante}
para aceptar un \omterm{def:b1:intermolecular-forces-solvents:hbond}{enlace de hidrógeno}.
\textbf{7.} El agua: polar, prótica y con la permitividad más alta; separa
y solvata los dos iones.
\textbf{8.} $78.4/1.89 = 41$: la atracción es 41 veces más débil en agua.
\textbf{9.} El acetonitrilo (o la propanona): lo bastante polar para
disolver la sal, y aprótico, para que el anión quede libre.
\textbf{10.} El hexano: inmiscible con el agua y apolar, como el diyodo. El
yodo pasa al hexano, la capa superior (el hexano es menos denso que el
agua).
\textbf{11.} El etanol es \omterm{def:b1:intermolecular-forces-solvents:miscibility}{miscible} con el agua: no se forma una segunda
capa.
\textbf{12.} Su oxígeno acepta \omterm{def:b1:intermolecular-forces-solvents:hbond}{enlaces de hidrógeno} del agua, pero no puede
dar ninguno y sus dos grupos etilo son \omterm{def:b1:intermolecular-forces-solvents:amphiphile}{hidrófobos}: solo se disuelve un
poco.
\textbf{13.} $-66.4 - (-85.0) = \qty{18.6}{\celsius}$ por \omterm{def:b1:periodicity:table}{periodo};
\ce{HF} extrapolado: $-85.0 - 18.6 = \qty{-103.6}{\celsius}$.
\textbf{14.} \ce{HF} hierve a \qty{19.6}{\celsius}, \qty{123}{\celsius}
por encima de la extrapolación.
\textbf{15.} $-62.5 - (-87.8) = \qty{25.3}{\celsius}$; \ce{NH3}
extrapolado: \qty{-113.1}{\celsius}; real, $-33.4$, \qty{80}{\celsius}
más.
\textbf{16.} $-88.1 - (-112.0) = \qty{23.9}{\celsius}$; \ce{CH4}
extrapolado: \qty{-135.9}{\celsius}; real, $-161.5$, por debajo. El metano
no forma \omterm{def:b1:intermolecular-forces-solvents:hbond}{enlaces de hidrógeno}; simplemente es pequeño y poco
polarizable.
\textbf{17.} Agua (unos \qty{179}{\celsius}) $>$ \ce{HF} (123) $>$
\ce{NH3} (80).
\textbf{18.} \ce{NH3} puede dar tres pero solo aceptar uno (un \omterm{def:b1:lewis-resonance-vsepr:lewis}{par no
enlazante}); \ce{HF} puede aceptar tres pero solo dar uno. El agua da dos
y acepta dos: cada molécula puede participar en cuatro \omterm{def:b1:intermolecular-forces-solvents:hbond}{enlaces de
hidrógeno}, una red tridimensional completa.
\textbf{19.} \ce{H2S} y \ce{H2Se} no forman \omterm{def:b1:intermolecular-forces-solvents:hbond}{enlaces de hidrógeno}
apreciables: sus temperaturas de ebullición siguen la tendencia de London
que seguiría el agua sin ellos. La propia agua es la anomalía que se
mide.
\textbf{20.} $-41.3 - (-60.3) = \qty{19.0}{\celsius}$ por \omterm{def:b1:periodicity:table}{periodo}.
\textbf{21.} $T(p) = -60.3 + 19.0\,(p - 3)$ en \unit{\celsius}.
\textbf{22.} $T(2) = -60.3 - 19.0 = \qty{-79.3}{\celsius}$.
\textbf{23.} $100.0 - (-79.3) = \qty{179}{\celsius}$.
\textbf{24.} Sin sus \omterm{def:b1:intermolecular-forces-solvents:hbond}{enlaces de hidrógeno}, el agua herviría hacia
\textbf{\qty{-79}{\celsius}} y sería un gas en la Tierra.
\end{solution}
